% ==========================================
% 3.2.a CONTEXT
% Responsible member: Bao
% ==========================================

\section*{3.2.a Context}
\addcontentsline{toc}{section}{3.2.a Context}

\subsection{Background}

In recent years, film photography has experienced a significant revival worldwide, including in Vietnam. Increasing numbers of photography enthusiasts, artists, and collectors are returning to analog photography because of its distinctive aesthetic qualities, creative experience, and the tangible value of physical film. Unlike digital photography, film photography requires a more deliberate shooting process and produces unique visual characteristics that cannot always be replicated by digital filters.

As the analog photography community continues to grow, the demand for supporting services such as film developing, film scanning, photo printing, film preservation, and digital archiving has also increased. However, the service ecosystem surrounding film photography remains highly fragmented. Photographers often need to search for Film Labs through social media platforms, online communities, messaging applications, or personal recommendations. Information regarding service availability, pricing, turnaround time, supported film formats, processing quality, and customer reviews is usually distributed across multiple channels.

\subsection{Current Problems}

Currently, there is no centralized digital platform that allows photographers to conveniently search, compare, book, and monitor film processing services. Customers usually communicate with Film Labs manually through social media, messaging applications, or phone calls to inquire about available services and order status. This communication process can be time-consuming and may lead to inconsistent information, delayed responses, and difficulties in tracking orders.

From the perspective of Film Labs, many laboratories still rely on manual workflows for receiving film rolls, recording customer information, monitoring processing stages, delivering digital scan files, and managing business operations. These fragmented processes may reduce operational efficiency and make it more difficult for Film Labs to scale their services as customer demand increases.

In addition, the analog photography community lacks an integrated digital ecosystem that combines professional services with community interaction. Photographers need not only reliable Film Lab services but also access to photography knowledge, equipment reviews, community discussions, workshops, photowalk events, and opportunities to buy or sell analog photography equipment. Digital scan files are also often stored separately across different cloud storage services, making it difficult for photographers to maintain a structured and searchable personal film archive.

\subsection{Project Motivation}

The growing popularity of film photography and the fragmented nature of the current service ecosystem create an opportunity for digital transformation. A centralized platform can improve the connection between photographers and Film Labs by providing transparent service information, standardized booking processes, real-time order tracking, and digital delivery of scanned photographs.

At the same time, the rapid development of Artificial Intelligence provides additional opportunities to improve the user experience. AI technologies can be applied to recommend suitable Film Labs based on user preferences, assist photographers in selecting appropriate film types, analyze scan quality, support intelligent search, and provide photography-related guidance through a domain-specific conversational assistant.

Therefore, this project proposes the development of an AI-powered digital platform that connects photographers, Film Labs, photography experts, delivery partners, and administrators within a unified ecosystem. The platform aims to digitalize film processing services, improve operational efficiency for Film Labs, facilitate knowledge sharing and community engagement, and provide intelligent and personalized services for analog photography enthusiasts.

\subsection{Project Objective and Significance}

The main objective of the project is to establish a scalable digital ecosystem that supports the complete lifecycle of analog photography. The proposed platform will allow photographers to discover Film Labs, compare available services, place service orders, request pickup and delivery, monitor processing progress, receive digital scans, and manage photographs through a personal Digital Film Archive.

For Film Labs, the platform will provide digital tools for service management, workflow monitoring, customer management, digital file delivery, and business analytics. Photography experts will be able to contribute tutorials, technical articles, equipment reviews, and professional knowledge to the community. Delivery partners can participate in the logistics process for film collection and delivery, while system administrators can monitor platform operations and ensure a reliable environment for all stakeholders.

By integrating Artificial Intelligence, cloud storage, digital asset management, and community-based services, the project is expected to modernize the analog photography ecosystem and create a sustainable platform that benefits both photographers and service providers.

